%! Equivalent Representations of Trigonometric Functions — Unit 3, Lesson 3.12 — Homework
\documentclass[10pt]{article}
\usepackage{precalculus-article}
\usepackage{precalculus-boxes}
\newcommand{\TallMath}[1]{$\displaystyle #1\rule[-1.4em]{0pt}{3.2em}$}

\begin{document}

\pageheader{Unit 3, Lesson 3.12}{Homework}
\namedateperiod

\vspace{0.10in}

\begin{notesbox}{Problems 1--6}

\textbf{1.}\quad Use a Pythagorean identity to simplify each expression to
a single term or constant.

\vspace{4pt}
\renewcommand{\arraystretch}{1.6}
\begin{tabular}{|l|c|}
\hline
\textbf{Expression} & \textbf{Simplified} \\
\hline
$5\sin^2\theta + 5\cos^2\theta$ & \\
\hline
$\tan^2\theta - \sec^2\theta$ & \\
\hline
$\sin^2\theta\cdot\csc^2\theta$ & \\
\hline
$\csc^2\theta - \cot^2\theta$ & \\
\hline
\end{tabular}

\vspace{0.14in}
\textbf{2.}\quad Use a sum or difference identity to find each exact value.

\vspace{4pt}
(a)\quad $\sin 105^\circ$
\hspace{0.6cm}(\textit{Write} $105^\circ = 60^\circ + 45^\circ$)

\vspace{4pt}
$\sin(60^\circ + 45^\circ)
= \sin 60^\circ \cos 45^\circ + \cos 60^\circ \sin 45^\circ
= $ \underline{\hspace{5cm}}

\vspace{8pt}
(b)\quad $\cos\!\left(\dfrac{7\pi}{12}\right)$
\hspace{0.6cm}(\textit{Hint:} $\dfrac{7\pi}{12} = \dfrac{\pi}{4} + \dfrac{\pi}{3}$)

\vspace{4pt}
$\cos\!\left(\dfrac{\pi}{4} + \dfrac{\pi}{3}\right)
= \cos\dfrac{\pi}{4}\cos\dfrac{\pi}{3} - \sin\dfrac{\pi}{4}\sin\dfrac{\pi}{3}
= $ \underline{\hspace{4cm}}

\vspace{0.14in}
\textbf{3.}\quad If $\sin\theta = \dfrac{5}{13}$ and $\theta$ is in Quadrant I, find $\cos 2\theta$
and $\sin 2\theta$. Show all work.

\vspace{4pt}
First find $\cos\theta$: $\cos\theta = $ \underline{\hspace{2.5cm}}

\vspace{4pt}
$\sin 2\theta = 2\sin\theta\cos\theta = $ \underline{\hspace{4.5cm}}

\vspace{4pt}
$\cos 2\theta = 1 - 2\sin^2\theta = $ \underline{\hspace{4.5cm}}

\vspace{0.14in}
\textbf{4.}\quad Verify the identity $\dfrac{1 - \cos 2\theta}{\sin 2\theta} = \tan\theta$.

\textit{Work from the left side only. Show each step.}

\writelines{5}

\vspace{0.14in}
\textbf{5.}\quad Solve $\sin^2\theta - \cos\theta - 1 = 0$ for $\theta \in [0, 2\pi)$.

\textit{(Hint: substitute $\sin^2\theta = 1 - \cos^2\theta$.)}

\writelines{5}

\textbf{Solutions:} $\theta \in \big\{ \underline{\hspace{4cm}} \big\}$

\end{notesbox}

\vspace{0.08in}

\begin{notesbox}{AP-Style Multiple Choice}

\textbf{6. (AP-Style MC)}\quad Which of the following is equivalent to
$\cos^2\theta - \sin^2\theta$ for all values of $\theta$?

\begin{enumerate}[label=(\Alph*), leftmargin=2em, itemsep=3pt]
  \item $1 - 2\sin^2\theta$
  \item $2\cos^2\theta$
  \item $2\sin\theta\cos\theta$
  \item $1 - \cos 2\theta$
  \item $\sin 2\theta$
\end{enumerate}

Answer: \underline{\hspace{1.5cm}}

\vspace{0.12in}
\textbf{7. (AP-Style MC)}\quad The equation $2\sin^2 x + 3\cos x - 3 = 0$ has
solutions in $[0, 2\pi)$. Which of the following is the complete solution set?

\begin{enumerate}[label=(\Alph*), leftmargin=2em, itemsep=3pt]
  \item $\left\{\dfrac{\pi}{3}\right\}$
  \item $\left\{0,\;\dfrac{\pi}{3},\;\dfrac{5\pi}{3}\right\}$
  \item $\left\{\dfrac{\pi}{3},\;\dfrac{5\pi}{3}\right\}$
  \item $\left\{0,\;\pi\right\}$
  \item $\left\{\dfrac{\pi}{3},\;\pi,\;\dfrac{5\pi}{3}\right\}$
\end{enumerate}

Answer: \underline{\hspace{1.5cm}}

\end{notesbox}

\vspace{0.08in}

\begin{tcolorbox}[colback=navylight, colframe=navy, arc=2mm,
    left=3mm, right=3mm, top=2mm, bottom=2mm]
\small\textbf{Preview --- Lesson 3.13: Trigonometric Equations and Inequalities}\\
You have solved trig equations on specific intervals like $[0, 2\pi)$.
In Lesson 3.13 we extend to general solutions of the form $\theta = \alpha + 2\pi k$
and solve trig \emph{inequalities}, interpreting solution sets on the number line.
\end{tcolorbox}

\end{document}
